\documentclass{ieeeaccess}
\usepackage{cite}
\usepackage{amsmath,amssymb,amsfonts}
\usepackage{algorithmic}
\usepackage{graphicx}
\usepackage{textcomp}
\def\BibTeX{{\rm B\kern-.05em{\sc i\kern-.025em b}\kern-.08em
    T\kern-.1667em\lower.7ex\hbox{E}\kern-.125emX}}
\begin{document}
\history{Date of publication xxxx 00, 0000, date of current version xxxx 00, 0000.}
\doi{10.1109/TQE.2025.DOI}
 
\title{Text Encoding in Quantum Natural Language Processing: A Systematic Survey}

\author{\uppercase{Danilo Cavaliere}\authorrefmark{1}, \uppercase{Giuseppe Fenza}\authorrefmark{1}, \uppercase{Lukasz Gajewski}\authorrefmark{1}, and \uppercase{Vincenzo Loia}\authorrefmark{1}}

\address[1]{Department of Management \& Innovation Systems, Università degli Studi di Salerno, Italy}

\markboth{TEXT ENCODING IN QUANTUM NATURAL LANGUAGE PROCESSING}{TEXT ENCODING IN QUANTUM NATURAL LANGUAGE PROCESSING}

\corresp{Corresponding authors: Danilo Cavaliere and Lukasz Gajewski (e-mails: \{dcavaliere, lgajewski\}@unisa.it).}

\begin{abstract}
% Ovviamente provvisorio, da raffinare con Risultati.
Quantum natural language processing (QNLP) applies quantum computational primitives (superposition, entanglement, and interference) to the representation and processing of linguistic information. Every QNLP system requires an encoding step, that is  a procedure by which discrete or distributed text representations are mapped onto qubit registers or parameterized quantum circuit inputs. This step governs model expressivity, qubit resource requirements, and trainability, hence it is rarely studied as an independent variable and no unified method has been already defined for analyzing and comapring different encodings. 
%unifying taxonomy has been designed yet. 
To fill the gap, this article presents a systematic literature review of text encoding strategies in QNLP and introduces a unifying taxonomy for analyzing and comparing them, covering 125 primary studies published between 2021 and 2026 and conducted under a protocol registered before the search was executed that fixes the search strategy, eligibility criteria, screening procedure, snowball sampling, and data extraction in advance. Six encoding families are identified and characterized: angle, amplitude, categorical grammar-based (e.g., DisCoCat), IQP/ZZ feature-map, data re-uploading, and hybrid classical-quantum embedding. Their distribution is analyzed across NLP tasks, simulation and hardware environments, qubit budgets, and performance relative to classical baselines. Noiseless quantum circuit simulation dominates the literature, real hardware execution remains rare, and claimed quantum advantage is predominantly conditional or absent. The encoding bottleneck, the dimensionality mismatch between linguistic representations and near-term qubit capacities, is formalized as a basis for resource-efficient QNLP algorithm and hardware design. The corpus and extracted metadata are released publicly to support reproducibility and future meta-analyses.


\end{abstract}

\begin{keywords}
Encoding bottleneck, parameterized quantum circuits, quantum advantage, quantum natural language processing, systematic literature review, text encoding.
\end{keywords}

\titlepgskip=-15pt

\maketitle

\section{Introduction}
\label{sec:introduction}
\PARstart{E}{very} quantum natural language processing (QNLP) system begins with the mandatory step of converting text into a quantum state. Before a single gate is applied or a single parameter is trained, a token, sentence, or document must be mapped onto qubits. This mapping (the encoding) determines how many qubits are consumed, how deep the resulting circuit must be, whether gradients can be computed efficiently, and whether any genuinely quantum property, such as entanglement, is brought to bear on the linguistic content. It is, in short, the most consequential design decision in the pipeline.

Despite this centrality, encoding is rarely the subject of dedicated study. Most QNLP papers dispatch it in a paragraph, introduce one strategy without justification, and move on to reporting accuracy numbers. The result, after five years of accelerating publication activity, is a fragmented body of work in which dozens of encoding families coexist but there is no common approach to analyze and compare them.  In fact, there is a strong lack of unified taxonomies, systematic comparisons, or shared understanding of encoding trade-offs in literature: existing QNLP surveys and reviews do not isolate encoding as a variable \cite{chen_survey_2024,pallavi_quantum_2025}, and the few explicitly comparative studies contrast model families or ans\"atze rather than encoding schemes \cite{bouakba_comparative_2025,castillo_comparative_2025}. This survey paper aims to fill the gap by conducting an extensive systematic literature review on topic and make text encoding in QNLP the primary subject of study.

\subsection{Scope}
This survey reviews primary research that (1) applies a genuine quantum circuit (simulated noiseless, simulated noisy, or executed on real hardware) to at least one NLP task, and (2) identifies, at minimum by name or type, the encoding scheme used to map text or text-derived features onto qubits or circuit parameters. This review follows a pre-registered SLR protocol (Section~\ref{sec:method}) adapted from established guidelines for secondary studies in computer science \cite{keele2007guidelines}.

Papers that use quantum-inspired or quantum-motivated classical algorithms without instantiating an actual parameterized quantum circuit are excluded. Other works excluded include survey and review papers, hardware-focused benchmarking studies that treat NLP datasets as a generic payload, and papers with insufficient methodological detail to assess the encoding scheme. The final corpus comprises 125 primary studies.

The scope is limited to the encoding step; quantum circuit design, ansatz expressibility, and quantum error correction are not surveyed except where they interact directly with encoding choices. Classical NLP methods are discussed only as baselines.

\subsection{Key Contributions}
This article makes \texttt{some (N)} contributions to the study of text encoding in QNLP: 
%Placeholder da compilare quando si hanno i contributi
\begin{enumerate}
\item[1)] A novel taxonomy to analyze and compare encoding methods.  
\item[2)] A systematic review of papers on NLP tasks using/proposing different quantum encodings.
\item[3)] A classification of the encoding techniques based on the taxonomy highlighting ....
\item[4)]
\item[5)]
\end{enumerate}

\subsection{Related Work}
Several prior works survey the QNLP literature, but they do not focus on text encoding as a crucial aspect to improve performance on NLP tasks. Varmantchaonala et al.\ \cite{varmantchaonala_quantum_2024} provide the most comprehensive QNLP survey to date, covering quantum algorithms, circuit architectures, and NLP task categories; encoding appears as an implementation detail rather than as an independent variable. Guarasci et al.\ \cite{guarasci_quantum_2022} offer an early but influential overview of QNLP challenges and opportunities, cataloging approaches from theoretical quantum-advantage arguments to hardware experiments, without isolating encoding as a dimension of comparison. Widdows et al.\ \cite{widdows_quantum_2024} examine word embeddings and sequential quantum models and discuss trade-offs between representation strategies, with more attention to encoding as a design choice than any other general survey. Bouakba and Belhadef \cite{bouakba_comparative_2025} compare compositional DisCoCat models against quantum neural network approaches experimentally; their taxonomy axis (compositional versus flat-vector encoding) overlaps with the one proposed in this survey, though the scope is limited to two model families. Huzzat et al.\ \cite{huzzat_primer_2026} is the most encoding-aware prior work: it explicitly addresses the mismatch between the dimensionality of text representations and the qubit capacity of NISQ classifiers, and discusses classical compression as a prerequisite for stable training. It is, however, a conceptual primer for a single model class rather than a systematic review across the literature.

This survey complements these works by treating encoding as the primary unit of analysis across 125 primary studies. Encoding family, qubit count, NLP task, execution environment, and quantum advantage claim were extracted for each study, yielding the first quantitative cross-study characterization of how text is encoded and at what resource cost.

\subsection{Structure}
The rest of the document is organized as follows: Section~\ref{sec:background} provides the background on quantum computing, compositional distributional semantics, and the role of encoding in QNLP. Section~\ref{sec:method} describes the systematic review protocol, including search strategy, inclusion and exclusion criteria. 
%Da continuare con altre sezioni.

\section{Background}
\label{sec:background}

\subsection{Quantum Computing Primitives}

A quantum bit (qubit) generalizes the classical bit by occupying a superposition of basis states $|0\rangle$ and $|1\rangle$:
\begin{equation}
|\psi\rangle = \alpha|0\rangle + \beta|1\rangle, \quad |\alpha|^2 + |\beta|^2 = 1,
\end{equation}
where $\alpha, \beta \in \mathbb{C}$ are probability amplitudes. A register of $n$ qubits occupies a $2^n$-dimensional Hilbert space $\mathcal{H} = \mathbb{C}^{2^n}$, and its state is a unit vector in that space. Entanglement between qubits allows correlations with no classical counterpart. Measurement collapses the state to a basis vector, destroying superposition and yielding a classical bitstring.

These properties, superposition, entanglement, and interference, are the computational primitives that QNLP seeks to bring to bear on language. Whether any of them confers a practical advantage over classical methods on NLP tasks is the central open question in the field and one of the dimensions analyzed in this survey.

\subsection{Parameterized Quantum Circuits}

The dominant computational model in near-term quantum machine learning is the parameterized quantum circuit (PQC), also called a variational quantum circuit \cite{cerezo2021variational}. A PQC applies a sequence of unitary gates to an initial state $|\psi_0\rangle$:
\begin{equation}
U(\boldsymbol{\theta}) = U_L(\boldsymbol{\theta}_L) \cdots U_1(\boldsymbol{\theta}_1),
\end{equation}
where $\boldsymbol{\theta} \in \mathbb{R}^p$ are trainable parameters. Measurement outcomes are summarized as expectation values $\langle O \rangle = \langle \psi_0 | U^\dagger(\boldsymbol{\theta})\, O\, U(\boldsymbol{\theta}) | \psi_0 \rangle$ that serve as model outputs. Gradients with respect to $\boldsymbol{\theta}$ are computed analytically via the parameter-shift rule \cite{mitarai2018quantum, schuld2019evaluating}, enabling optimization by standard gradient-based methods. PQCs are the circuit model underlying five of the six encoding families surveyed here; the exception is grammatical encoding, which derives circuit structure from sentence syntax rather than from a fixed ansatz.

PQCs are constrained by the limits of noisy intermediate-scale quantum (NISQ) hardware \cite{preskill2018quantum}: shallow circuits, limited qubit counts, and significant gate error rates. These constraints bear directly on the encoding step, since every additional qubit or gate layer required to load a text representation consumes resources that could otherwise be used for the learning task.

\subsection{Text Encoding}

Text encoding is the procedure that maps a text input $x$, a token, sentence, or document, to a quantum state $|\phi(x)\rangle \in \mathcal{H}$ or to a parameter vector $\boldsymbol{\theta}(x)$ that initializes a PQC. It is the mandatory interface between classical language data and the quantum circuit. The input $x$ may be a raw token sequence, a classical embedding vector from a pre-trained model, or a parse tree carrying grammatical structure; which representation is used, and how it is loaded onto qubits, defines the encoding family.

Encoding choices have cascading consequences. The number of qubits required, the depth of the state-preparation circuit, the presence or absence of linguistically meaningful entanglement, and the trainability of the resulting model all depend on this single decision. Despite this, most QNLP papers treat encoding as a fixed implementation detail rather than as a variable to be studied.

\subsection{Compositional Distributional Semantics and DisCoCat}

One of the two main traditions in QNLP originates in categorical linguistics. Coecke, Sadrzadeh, and Clark \cite{coecke2010mathematical} showed that the compact closed categorical structure of pregroup grammars is the same structure that governs quantum processes in Hilbert space. In a pregroup grammar, each word is assigned a grammatical type, and sentences are parsed by reducing compound types through tensor contractions. Because these reductions obey the same diagrammatic calculus as quantum teleportation steps, a grammatical parse tree translates directly into a quantum circuit: words become parameterized quantum states, and grammatical reductions become cups in a string diagram.

This is the DisCoCat (Distributional Compositional Categorical) model, operationalized as an end-to-end pipeline by \cite{lorenz_qnlp_2023} and executed on near-term hardware in \cite{meichanetzidis_quantum_2021, meichanetzidis_grammar-aware_2023} via the \textsc{lambeq} toolkit. The circuit structure in DisCoCat is not fixed in advance: it is determined by the parse tree, so syntactically distinct sentences produce structurally distinct circuits. This property makes the grammatical encoding family categorically different from all other families, which treat the sentence as a flat feature vector.

\subsection{Quantum Machine Learning for Text}

The second tradition approaches NLP as a classification problem and applies quantum feature maps \cite{schuld2019quantum} and variational circuits \cite{cerezo2021variational} to text-derived feature vectors. Classical embeddings, word2vec, GloVe, or transformer outputs, are compressed to match available qubit counts and loaded into a PQC via rotation gates or amplitude encoding. The circuit is then trained end-to-end on labeled examples. Because this approach makes no assumptions about grammatical structure, it applies to arbitrary text inputs without a parser and scales more easily to realistic dataset sizes. Hybrid quantum--classical architectures, in which a large classical language model handles high-dimensional representation and a small quantum circuit acts as a downstream head, represent the logical endpoint of this design philosophy and have grown in prevalence since 2023.

The two traditions differ not only in how they encode text but in what they claim the quantum component contributes: compositional models argue for structural alignment between grammar and quantum computation; QML models argue for expressivity or kernel separability. 

\section{Review Methodology}
\label{sec:method}

\Figure[t!](topskip=0pt, botskip=0pt, midskip=0pt)[width=\textwidth]{PRISMA_Chart.jpg}{PRISMA flow diagram. Final corpus: $n = 125$ (95 from database search, 30 from snowball sampling).\label{fig:prisma}}

This systematic literature review follows the guidelines for secondary studies in software engineering established by Keele and Charters \cite{keele2007guidelines}. The protocol was defined before the search was executed and covers the search strategy, eligibility criteria, screening procedure, snowball sampling, and data extraction process.

\subsection{Search Strategy}
Two complementary search queries were executed on two electronic databases (Scopus and IEEE Xplore) to provide broad coverage of both the quantum machine learning and the compositional QNLP literature. The first query targets the intersection of quantum computing or machine learning with NLP and encoding-related terms; the second targets QNLP-specific terminology including the DisCoCat and DisCoCirc frameworks, established pipelines that implement them such as the \textsc{lambeq} toolkit, and categorical linguistic concepts. The searches were restricted to English-language journal articles and conference papers published between 2021 and 2026, ends included, but excluding editorials, book chapters, and erratum notices. The searches returned 269 results from Scopus and 217 from IEEE Xplore, for a combined total of 486 records. After removing 110 duplicate records, 376 unique records were obtained.

\subsection{Eligibility Criteria}
\label{sec:criteria}

For paper selection, several criteria have been considered. Then, a document has been selected if and only if it met all the six inclusion eligibility criteria listed below simultaneously.  

\begin{itemize}
    \item \textit{IC1: NLP task grounding.} The method proposed in the paper is applied to, or evaluated on, at least one natural language processing task such as classification, similarity, sequence labeling, question answering, or generation.

    \item \textit{IC2: Genuine quantum computation.} The method is implemented and evaluated via at least a noiseless quantum circuit simulation (exact statevector simulation or equivalent) or executed on real quantum hardware. Classical algorithms that are merely quantum-inspired, such as tensor-network models or classical approximations of quantum kernels, do not satisfy this criterion.

    \item \textit{IC3: Encoding scheme identified.} The paper names or identifies at least the type of encoding used to map text or text-derived features onto qubits, quantum states, or circuit parameters. A named or typed encoding mention is sufficient; circuit diagrams or full mathematical derivations are not required.

    \item \textit{IC4: Peer-reviewed or preprint venue.} Published in a peer-reviewed journal, conference, or workshop, or as a preprint with sufficient methodological detail to assess IC1--IC3.

    \item \textit{IC5: Language and access.} Full text available in English and accessible for review.

    \item \textit{IC6: Time window.} Published between 2021 and 2026 inclusive.
\end{itemize}

Even though a paper satisfied all the criteria listed above, it has been excluded if one of the following criteria was satisfied: 
\begin{itemize}
    \item \textit{EC1: No NLP application.} Quantum ML papers applied exclusively to non-text domains with no transferable text-encoding discussion.

    \item \textit{EC2: Classical-only encoding.} Papers using encoding purely in the classical sense (tokenization, positional encoding, one-hot) with no quantum circuit involved.

    \item \textit{EC3: Insufficient detail.} Abstracts, extended abstracts, or short papers that name an encoding but provide no formal description or reproducible detail.

    \item \textit{EC4: Duplicate or superseded.} Conference version superseded by an extended journal version; only the more complete version is retained.

    \item \textit{EC5: Pure hardware focus.} Papers focused solely on hardware benchmarking that use an NLP dataset only as a generic payload with no discussion of linguistic encoding properties.

    \item \textit{EC6: Quantum-inspired, classical surrogate, or unverifiable simulation.} Papers that use quantum-derived mathematical formalism without instantiating an actual parameterized quantum circuit, report only classical surrogates of quantum behavior, or claim quantum simulation without providing evidence of framework or circuit use.

    \item \textit{EC7: Survey or review paper.} Papers that are primarily secondary studies rather than original research contributions; these were retained as reference material and are discussed in Section~\ref{sec:introduction}.
\end{itemize}

\subsection{Screening and Selection}
\label{sec:screening}

Screening proceeded in two stages. In the first stage, all 376 records were assessed on title and abstract alone for broad topical relevance to quantum computing applied to text. No IC/EC criteria were applied at this stage; the goal was only to remove records that were clearly outside the scope of the review. This eliminated 234 records, leaving 142 reports sought for retrieval, of which 6 could not be retrieved, leaving 136 reports assessed for eligibility.

In the second stage, the full text of each of the 136 reports was assessed against the complete set of inclusion and exclusion criteria defined in Section~\ref{sec:criteria}. This resulted in 95 papers included in the primary corpus. The remaining 41 retrieved reports were excluded or retained as surveys, mainly under IC2 (no genuine quantum circuit; $n = 15$), IC3 (encoding scheme not identified; $n = 6$), EC6 (quantum-inspired, classical surrogate, or unverifiable simulation; $n = 3$), and EC7 (survey or review; $n = 11$); six records were excluded under other criteria.

\subsection{Snowball Sampling}
To reduce the risk of missing relevant work not indexed by the two databases, two rounds of backward and forward snowball sampling \cite{wohlin2014guidelines} were applied to the surveys excluded under EC7. Reference lists and citing papers yielded 1222 unique candidate records. In the first stage, these were screened on title and abstract for broad topical relevance. In this way, the set was reduced to 242 reports sought for retrieval, of which 37 could not be retrieved, leaving 205 reports assessed for eligibility. In the second stage, the full text of each report was assessed against the complete IC/EC protocol. This resulted in 30 additional papers included. The remaining 175 retrieved reports were excluded or retained as surveys, mainly under IC1 (no NLP task; $n = 74$), IC2 (no genuine quantum circuit; $n = 52$), IC6 (outside the time window; $n = 20$), IC3 (encoding scheme not identified; $n = 4$), and EC7 (survey or review; $n = 22$); three records were excluded under other criteria.

\subsection{Final Corpus}
The final primary corpus comprises 125 papers: 95 from database search and 30 from snowball sampling. Fig.~\ref{fig:prisma} summarizes the selection process following the PRISMA flow diagram convention.

\section{Taxonomy of Text Encoding Strategies}
\label{sec:taxonomy}

The encoding problem in QNLP is not a single operation. Depending on the study, text can enter the quantum computation as rotation parameters, as amplitudes of a prepared state, as a grammar-derived circuit, as a structured feature map, or through a more specialized interface such as repeated data injection, a hybrid neural representation, or an observable. These mechanisms are often discussed under the same term, ``encoding,'' although they make different assumptions about where linguistic information is represented and which quantum resource is being used.

This section shows how these mechanisms are organized into the proposed taxonomy as shown in Table~\ref{tab:encoding_taxonomy}. The classification is based on the point at which text or text-derived features enter the quantum computation. We therefore do not classify a study by the optimizer, measurement rule, downstream classifier, or ansatz alone. These elements are treated as design choices that can be combined with an encoding category. This distinction is important because an applied QNLP pipeline is usually a combination of several decisions: a classical text representation, a quantum encoding interface, a trainable or fixed circuit, and a measurement or kernel rule.

The final analyzed corpus contains 125 primary studies. Four major encoding families are observed: angle encoding ($n=51$), amplitude encoding ($n=28$), syntactic/DisCoCat encoding ($n=29$), and structured feature-map encoding ($n=9$). A smaller residual group ($n=8$) contains specialized interfaces: data re-uploading ($n=4$), hybrid classical--quantum embedding ($n=3$), and Hamiltonian/observable encoding ($n=1$). To use the taxonomy, one may first identify whether the text representation is loaded as a flat vector, imposed through linguistic structure, or used to define a kernel geometry. The next design question is whether the encoding is intended to minimize qubit count, preserve linguistic composition, define a kernel geometry, or improve trainable circuit expressivity. Not all combinations are equally practical on near-term devices. The main tradeoff is between width, depth, state-preparation cost, and the amount of linguistic structure preserved by the quantum circuit \cite{preskill2018quantum,cerezo2021variational,ruskanda_evaluating_2025}.

\begin{table*}[t]
\caption{Taxonomy of text encoding strategies in the corrected primary corpus.}
\label{tab:encoding_taxonomy}
\centering
\scriptsize
\begin{tabular}{@{}p{0.18\textwidth} p{0.09\textwidth} p{0.45\textwidth} p{0.20\textwidth}@{}}
\hline
\textbf{Encoding family} & \textbf{Corpus count} & \textbf{Core mechanism} & \textbf{Typical model families} \\
\hline
Angle encoding & 51 & Text features are mapped to single-qubit rotation parameters, typically through $R_x$, $R_y$, or $R_z$ gates. & VQC, hybrid neural, recurrent neural, quantum transformer \\
Amplitude encoding & 28 & A normalized feature vector is loaded into the amplitudes of an $n$-qubit state in a $2^n$-dimensional Hilbert space. & VQC, quantum kernel, quantum transformer, hybrid neural \\
Syntactic/DisCoCat encoding & 29 & Sentence grammar determines a string diagram or circuit; word meanings are attached to grammatical types and composed by reductions. & Compositional QNLP \\
Structured feature maps & 9 & Classical text features parameterize structured circuits such as IQP, ZZFeatureMap, ZFeatureMap, PauliFeatureMap, or Pauli-Y variants. & Quantum kernel, VQC, hybrid neural, quantum transformer \\
Data re-uploading & 4 & Input features are injected repeatedly across circuit layers, alternating data-dependent rotations with trainable transformations. & VQC, recurrent neural \\
Hybrid embedding & 3 & A classical encoder first produces a compact text representation that is then passed to a quantum layer. & Hybrid neural, VQC \\
Hamiltonian/observable encoding & 1 & Text-derived information is represented through Pauli-string Hamiltonians or observables rather than only through prepared quantum states. & VQC \\
\hline
\end{tabular}
\end{table*}

\subsection{Angle Encoding}
\label{subsec:angle_encoding}

Angle encoding is the dominant category in the corpus. In this scheme, a classical vector $\mathbf{x} = (x_1,\ldots,x_d)$ is mapped to rotation parameters,
\begin{equation}
|\phi(\mathbf{x})\rangle =
\prod_{j=1}^{d'} R_{\alpha_j}(s_j(\mathbf{x})) |0\rangle^{\otimes n},
\end{equation}
where $R_{\alpha_j}$ is a rotation around the $x$, $y$, or $z$ axis and $s_j(\cdot)$ denotes scaling, projection, or feature selection. The input is rarely raw text. It is usually a TF-IDF vector, a word or sentence embedding, a transformer representation, or a compact feature vector produced by a preceding classical model.

The prevalence of angle encoding follows from its implementation simplicity. Rotations fit naturally into parameterized quantum circuits and support the standard gradient-based training workflow of variational algorithms, and they avoid arbitrary-state preparation, which makes them convenient for small simulated or hardware-oriented experiments \cite{mitarai2018quantum,schuld2019evaluating}. What makes the family a single encoding category is this shared rotation mechanism; within it, the corpus studies differ along two encoding-level axes: the rotation construction itself and the classical source that supplies the angles.

The first axis is the rotation construction. The most common pattern assigns one feature to one single-axis rotation, typically an $R_Y$ rotation as in the default angle-embedding routine \cite{ruskanda_evaluating_2025,kesiku_quantum_2025,diaz-ortiz_strongly_2024}, though $R_X$ is also used \cite{day_qnet_2023,zhou_new_2022,tiwari_quantum_2024} and some studies apply combined $R_Y$ and $R_Z$ rotations per qubit \cite{stein_applying_2023}. A denser variant stores more than one feature on each qubit through successive rotations about different axes, trading a shorter register for greater loading depth \cite{kesiku_unlocking_2026}.

The second axis is the classical source of the angles, and it is where the studies vary most, because the rotation parameters are numbers rather than text and so the encoding's content is fixed by the front-end that produces them. Some studies load sparse lexical statistics such as TF-IDF, bag-of-words, or $n$-gram vectors \cite{boopalamani_quantum-enhanced_2025,payares_quantum_2023,kabir_qtopic_2026}; others load static word embeddings \cite{di_sipio_dawn_2022,diaz-ortiz_strongly_2024,ohno_quantum_2025}; and a large, growing group compresses contextual transformer embeddings, for instance reducing a 768-dimensional BERT or DistilBERT vector to a handful of features before rotation \cite{ahmad_hybrid_2023,kumar_hybrid_2026,suleiman_hdqe_2026,yu_application_2024}. A further group derives the angles from deeper classical feature extractors or from handcrafted linguistic features \cite{rezaei_hybrid_2026,mukesh_qvila_2024,mompelat_compact_2026}. In each case a dimensionality-reduction step, whether PCA, a linear projection, or an autoencoder, sits between the text representation and the rotations and is effectively part of the encoding, since it determines what survives into the quantum circuit.

The limitation is inherent in this design. Encoding cost grows with the number of retained features, so a model must spend more rotations, more qubits, or more layers as it keeps more of the input. Angle encoding therefore pushes the hard part of the problem onto classical dimensionality reduction, which makes it practical for NISQ-scale prototypes but also means the quantum circuit often operates on a heavily compressed view of the linguistic input.

\subsection{Amplitude Encoding}
\label{subsec:amplitude_encoding}

Amplitude encoding maps a normalized vector $\mathbf{x} \in \mathbb{R}^{2^n}$ or $\mathbb{C}^{2^n}$ to a quantum state,
\begin{equation}
|\phi(\mathbf{x})\rangle = \sum_{i=0}^{2^n-1} x_i |i\rangle,
\quad \sum_i |x_i|^2 = 1.
\end{equation}
For NLP, the appeal is immediate. A high-dimensional sentence embedding can, in principle, be represented with logarithmically many qubits: a 768-dimensional transformer vector, for example, can be padded and normalized into a ten-qubit state. This density is why amplitude encoding is most often paired with dense contextual representations, above all BERT, ELECTRA, and other transformer embeddings, and to a lesser extent with static word embeddings or TF-IDF and bag-of-words vectors \cite{buonaiuto_quantum_2024,guarasci_applying_2025,cavar_old_2026,zeguendry_quantum_2024}.

Because the register holds exactly $2^n$ amplitudes, the encoding is defined as much by how the source vector is fitted to that dimension as by the loading step itself. When the embedding is already close to a power of two, it is simply normalized and zero-padded, as when a 768-dimensional vector is loaded into ten qubits \cite{guarasci_applying_2025,buonaiuto_multilingual_2025}. When it is not, or when the qubit budget is smaller, a classical reduction such as PCA, t-SNE, or UMAP is applied first, so that the reduction and the amplitude map together fix the geometry the circuit actually sees \cite{korrapati_quantum_2025,zhang_entangled_2024,tufchi_quantum-enhanced_2026}.

The reason to accept this encoding is representational. Storing a large vector in a small register makes overlap- and inner-product-based operations natural, because the normalized state \emph{is} the embedding. SWAP-test similarity and quantum-kernel pipelines exploit this directly, letting the geometry of the encoded state stand in for the similarity structure of the embedding space \cite{cavar_word_2025,zeguendry_quantum_2024,schuld2019quantum}. A distinct thread goes further and encodes complex-valued rather than real-valued amplitudes, representing tokens as complex state vectors and building complex-valued semantic operations on top of them, which uses a property of the amplitude space that has no direct classical-vector analogue \cite{routh_quantum_2026,cavar_old_2026}.

The signature difficulty is state preparation, and it is where the studies differ most in how honestly the cost is treated. Most rely on exact library routines that prepare an arbitrary amplitude vector, with the preparation depth absorbed into a simulator call \cite{buonaiuto_quantum_2024,guarasci_applying_2025}. A smaller set confronts the cost directly through divide-and-conquer amplitude encoding, which trades circuit depth against ancilla width \cite{zhang_entangled_2024}. Loading an arbitrary vector generally requires a circuit whose depth can offset the qubit savings, so experiments that hide this behind high-level routines remain useful for comparing representation strategies under a fixed qubit budget but should not be read as hardware-ready evidence unless the preparation cost is included.

\subsection{Syntactic and DisCoCat Encoding}
\label{subsec:syntactic_encoding}

Syntactic encoding is the characteristic interface of compositional QNLP. Instead of first reducing a sentence to a flat feature vector, the model begins with a grammatical derivation. Words are assigned grammatical types, the sentence is reduced according to a pregroup or related grammar, and the resulting string diagram is translated into a tensor network or quantum circuit \cite{coecke2010mathematical}. What makes the family a single encoding category is a shared two-stage construction: a diagram that fixes the circuit \emph{topology} from grammar, followed by a rule that turns that diagram into a parameterized qubit circuit. The 29 studies in the corpus differ mainly along these two stages, and it is these differences, rather than the downstream task, that define the encoding-level variants below.

The first stage is how the diagram is obtained. The canonical route is a pregroup derivation, in which each word receives a pregroup type and grammatical reductions become the wires that connect word boxes; pregroup parsing is efficient and underlies most DisCoCat implementations \cite{meichanetzidis_quantum_2021,meichanetzidis_grammar-aware_2023,lorenz_qnlp_2023}. In practice this parse is often produced by a broader-coverage CCG parser and then rewritten into DisCoCat form, as in studies that derive their diagrams with the Bobcat or DepCCG parsers before compilation \cite{abiwardani_generalized_2025,ruskanda_quantum_2022,r_quantum_2023}. A second variant replaces the parse with a fixed \emph{reader} that imposes a uniform composition rule instead of syntax: a bag-of-words reader composes words through a single commutative spider, and a word-sequence reader composes them left to right into a tensor train, both of which appear as structure-free counterparts to the grammar-based diagrams \cite{lorenz_qnlp_2023,liu_research_2024,peral-garcia_comparing_2024}. These readers trade syntactic fidelity for shallower, more uniform circuits and are frequently used as structure-free reference points within DisCoCat studies. A third variant changes the scope of the diagram itself, lifting the sentence-level construction to the discourse level so that pronouns and cross-sentence reference are represented, as in DisCoCirc-style coreference models \cite{wazni_large_2024}.

The second stage is how the diagram becomes a qubit circuit. Each atomic type is assigned a fixed number of qubits, so the noun and sentence types set the circuit width, grammatical reductions compile to Bell-style effects between wires, and function words such as relative or possessive pronouns are realized as Kronecker or GHZ spiders \cite{meichanetzidis_grammar-aware_2023}. Word boxes are then filled by a parameterized ansatz, and the dominant choices are the IQP ansatz used by default in \textsc{lambeq} and, where circuit expressivity is compared, the Sim-14, Sim-15, and strongly entangling ans\"atze \cite{lorenz_qnlp_2023,castillo_comparative_2025}. Some studies design task-specific ans\"atze on top of the same grammatical skeleton rather than proposing a new diagram \cite{ruskanda_simple_2023,abiwardani_generalized_2025}. A final variant intervenes between the two stages by rewriting the diagram before compilation: ZX-calculus is used to simplify the grammar-derived circuit and reduce its resource cost \cite{sreedhar_towards_2025}.

This category differs from vector-loading schemes. The circuit is not only a host for numerical features; its wiring reflects syntactic composition. This gives syntactic encoding a clear interpretability advantage. Grammatical roles, reductions, and word interactions remain visible in the model structure, which explains why nearly all compositional papers in the corpus use this category.

The tradeoff is flexibility. Syntactic encoding depends on parsers, grammar formalisms, and sentence structures that can be translated into manageable circuits. Many experiments therefore use short, controlled, or grammar-friendly datasets. As sentences become longer, noisier, multilingual, or document-level, the parser and circuit-growth requirements become more restrictive. Thus, syntactic encoding is the most linguistically structured category, but it is also the category that most depends on the availability of a reliable grammar-to-circuit pipeline.

\subsection{Structured Feature-Map Encoding}
\label{subsec:iqp_encoding}

Structured feature maps encode data through circuits with phase rotations, Pauli operators, and entangling gates whose parameters depend on products or interactions of input features, so that the encoded state defines an implicit feature space in which examples are compared by state similarity \cite{schuld2019quantum}. Within the corpus the map is drawn from a small set of families. The most common is the ZZ feature map, which introduces pairwise feature interactions through $ZZ$ entangling rotations \cite{widdows_near-term_2024,baik_neural_2025,sunkavalli_experimentation_2024}. IQP embeddings, which apply diagonal phase gates in the Hadamard basis, are used similarly \cite{djemmal_arabert-qc_2025,ruskanda_quantum-enhanced_2023}. A few studies construct higher-order maps from selected Pauli terms, for example combined $X$, $Y$, and $ZZ$ gates or a hierarchy of Pauli-$Y$ and second-order $YY$ interactions, in order to encode richer correlations \cite{katyayan_supervised_2023,ruskanda_quantum-enhanced_2023}. One study instead reads the feature map out through local observables, using a projected quantum kernel to obtain a lower-variance similarity estimate \cite{keith-norambuena_quantum_2026}.

Because both the qubit count and the number of encoded interactions grow with the input dimension, a structured feature map is almost always preceded by an aggressive classical reduction. In practice this is usually PCA to on the order of ten to a few dozen components, and occasionally a wavelet or related projection, applied to an already compressed text representation such as a TF-IDF vector or a reduced transformer embedding \cite{mohammad_masum_hybrid_2023,djemmal_arabert-qc_2025,jagatheesan_quantum_2026}. The reduction is thus part of the encoding, since it fixes both the qubit budget and which feature interactions the map can express.

The benefit is that feature interactions are built into the map. Rather than using the circuit only to host features, the feature map defines the geometry in which examples are compared, which makes the approach attractive for testing whether a quantum kernel yields a useful similarity space for text.

The limitation is that the feature map becomes a strong modeling assumption. Performance depends on whether the chosen interactions are relevant to the linguistic task and on the preceding compression, so a structured feature map offers a principled kernel-oriented interface whose usefulness is set jointly by map design and the reduction that feeds it.

\subsection{Other Specialized Encoding Interfaces}
\label{subsec:other_encoding_interfaces}

The remaining studies use mechanisms that are important but too sparse or heterogeneous to treat as major standalone families. Data re-uploading addresses the small-register setting by injecting the input features multiple times across the depth of a circuit: rather than loading the data once and then applying only trainable gates, the model alternates data-dependent rotations with trainable transformations, so expressivity grows with depth rather than with qubit count. The four corpus studies differ mainly in what is re-uploaded and where. One implements a bricklayer of dense rotation gates in which repeated feature injection, rather than one-shot angle loading, is the defining encoding choice \cite{hamilton_characterizing_2023}. A recurrent variant re-uploads code-token embeddings across successive quantum embedding layers, so the same feature stream re-enters the circuit at every recurrent step \cite{song_recurrent_2024}. A third injects lexical features repeatedly into a seven-qubit NISQ classifier, explicitly trading qubit width for circuit depth \cite{silver_lexiql_2024}, and a fourth re-uploads sentence-transformer embeddings across the layers of an eighteen-qubit circuit, using qubit count and depth jointly to control capacity in a quantum LLM fine-tuning setting \cite{kim_quantum_2025}. The shared advantage is that expressivity can be raised without more qubits; the shared drawback is that the encoding burden moves from width to depth.

Hybrid embedding is a system-level interface rather than a single gate-level rule. A classical model first produces a compact representation, and a quantum layer then consumes it through an angle, amplitude, or less explicitly specified feature-injection mechanism. In the corpus the three studies are distinguished mainly by how aggressively the classical stage compresses and how clearly the quantum loading is described. One compresses fused text and image embeddings into an eight-dimensional latent vector and matches it to an eight-qubit VQC for emotion recognition in dialogue \cite{castrignano_toward_2025}. Another compresses GloVe and fastText code-token embeddings into a minimal intermediate representation before a four-qubit vulnerability classifier, without specifying the gate-level loading rule \cite{akter_automated_2023}, and a third fuses GloVe text features with VGG16 visual features through a quantum layer for misinformation detection, again with no stated state-preparation method \cite{bikku_quantum-enhanced_2025}. This underspecification is characteristic of the family and feeds its main difficulty, attribution: when most semantic processing happens in the classical stage, the independent contribution of the quantum layer is hard to isolate.

Hamiltonian and observable encoding is even more specialized. The single corpus study maps text-derived vectors onto Pauli-string Hamiltonians and reads them out as quantum observables rather than preparing them as amplitude or rotation states, which avoids conventional state encoding entirely and is reported to keep qubit and gate complexity logarithmic \cite{tiblias_efficient_2026}. This relocates the loading problem from state preparation to observable design and may sidestep some state-preparation costs, but with only one study in the corpus the evidence base is too small to support a separate taxonomy family.

\subsection{Taxonomic Implications}
\label{subsec:taxonomy_implications}

The taxonomy makes clear that encoding is not a neutral preprocessing choice in QNLP. It determines which linguistic information reaches the quantum circuit, how much classical compression is required, and where the main resource bottleneck appears. Angle and amplitude encodings start from numerical text representations and differ mainly in how they trade circuit width, depth, and state-preparation cost. Syntactic encoding instead gives the grammatical analysis a structural role in the circuit. Structured feature maps define the comparison space used by kernel methods, while the specialized interfaces modify these patterns through repeated data injection, classical--quantum embedding layers, or observable-based representations.

This distinction matters because many cross-paper comparisons implicitly compare different views of the input rather than only different learning architectures. A VQC using compressed word or sentence vectors, a DisCoCat circuit generated from a parser, and a kernel model built around a designed feature map make different commitments about what should be preserved from the original text. Their results are therefore shaped jointly by the classical representation, the compression method, the encoding circuit, the ansatz, the measurement rule, and the execution backend.

The taxonomy is therefore used in the remainder of the survey to organize the evidence, not to rank the encodings. The available studies do not yet show that any single family is generally superior across QNLP tasks. Instead, they show that the encoding stage frames the experiment itself: it defines the form of the linguistic input, constrains the feasible circuit resources, and conditions how performance, scalability, and quantum-advantage claims should be read.

\section{Comparative Analysis of Encoding Families}
\label{sec:comparative_analysis}

This section uses the taxonomy to assess how far the main encoding families have developed in practice. Rather than revisiting the definition of each family, the focus is on what the corpus reveals about their use: which encodings have become standard design choices, which remain exploratory, and how strongly the reported results are supported by datasets, baselines, hardware assumptions, and scalability analysis.

The analysis is based on the corrected corpus of 125 primary studies. Each study is assigned to its dominant encoding interface, although most implementations combine several steps before measurement: text representation, feature reduction, circuit encoding, a trainable or fixed quantum model, and a readout rule. This matters because an observed result is rarely attributable to the encoding alone. It often reflects the full pipeline in which the encoding is embedded.

\Figure[t!](topskip=0pt, botskip=0pt, midskip=0pt)[width=\textwidth]{figures/encoding_by_model_family.pdf}{Encoding family by model family in the final primary corpus. Cell values report the number of studies assigned to each encoding--model-family pair.\label{fig:encoding_model_family}}

Before examining the families individually, it is worth noting that the choice of encoding is not independent of the model family in which it appears. Figure~\ref{fig:encoding_model_family} crosses the two dimensions, and the resulting distribution is far from uniform: most encodings concentrate in one or two model families rather than spreading evenly across them. The three dominant encodings account for 108 of the 125 studies, yet two of them are all but fixed by the architecture they serve. Syntactic/DisCoCat encoding and the compositional model family form a two-way lock, since all 29 syntactic studies are compositional and all 29 compositional studies are syntactic, and a weaker version of the same coupling ties structured feature maps to quantum kernels, where six of the nine structured-map studies live. Angle encoding, by contrast, is the general-purpose default shared across the remaining gate-model architectures, and amplitude encoding co-resides with it inside VQCs and kernel methods. This is why comparing ``encoding families'' across the corpus is partly comparing the model families they are bound to, and it motivates the approach of the four subsections that follow, which examine each established family in turn by holding the encoding fixed and asking how its use changes across the contexts in which it appears; the specialized interfaces are treated last.

\subsection{Angle Encoding Across Model Families}
\label{subsec:angle_across_families}

Angle encoding is the one interface that is not bound to a single architecture, so it is the natural place to begin: rather than asking which encoding each model family prefers, we hold the encoding fixed and ask how the same data-dependent rotation is repurposed as the surrounding model changes. The 51 angle-encoding studies distribute across five model families, VQCs (20), hybrid neural models (10), recurrent models (10), quantum transformers (9), and quantum kernels (2), and their treatment of the rotation differs enough that ``angle encoding'' names a role rather than a fixed operation. Table~\ref{tab:angle_by_family} summarizes the split; the remainder of this subsection reads it family by family. Qubit budgets are reported by 34 of the 51 studies and, with one structural exception, stay within a narrow $2$--$8$ band, so the differences below are differences in \emph{where} the rotation sits, not in how wide a register it drives.

\begin{table*}[t]
\caption{Angle encoding split by model family in the final primary corpus.}
\label{tab:angle_by_family}
\centering
\scriptsize
\begin{tabular}{@{}p{0.15\textwidth} p{0.06\textwidth} p{0.13\textwidth} p{0.26\textwidth} p{0.27\textwidth}@{}}
\hline
\textbf{Model family} & \textbf{Studies} & \textbf{Median qubits (range)} & \textbf{Typical classical front-end} & \textbf{Role of the angle rotation} \\
\hline
VQC & 20 & 4.5 (2--16) & TF-IDF, BERT/SBERT/DistilBERT, handcrafted features & Sole quantum front door before a trainable ansatz \\
Hybrid neural & 10 & 4.5 (2--8) & Output of a classical transformer/RNN encoder & Submodule consuming an already-compact representation \\
Recurrent neural & 10 & 4 (4--8) & Word or BERT embeddings, per token & Data input inside the recurrent cell, re-entered each step \\
Quantum transformer & 9 & 4 (4--28) & Word embeddings, multimodal features & State preparation for quantum attention \\
Quantum kernel & 2 & 4 (4) & TF-IDF or compressed feature vectors & Definition of a similarity-circuit geometry \\
\hline
\end{tabular}
\end{table*}

In VQCs, angle encoding appears in its canonical form: the rotation layer \emph{is} the quantum front door, followed only by a trainable ansatz and a measurement, so the encoding carries the entire burden of moving text into the circuit. This is also the family with the widest spread of classical front-ends, which is why the compression the encoding demands is most visible here. Studies load TF-IDF or bag-of-words features \cite{thota_comparison_2025}, transformer embeddings reduced by one or two orders of magnitude, from BERT's 768 dimensions to ten \cite{ahmad_hybrid_2023}, from DistilBERT to four \cite{kumar_hybrid_2026}, or as far down as two values before an $R_x$ layer \cite{altintas_beyond_2025}, and in a few cases handcrafted linguistic descriptors such as character entropy and diacritic counts \cite{mompelat_compact_2026}. Because the circuit is the model rather than a component of it, the VQC setting also hosts almost all of the corpus's controlled encoding comparisons: a QCNN study finds angle encoding the best of angle, amplitude, and IQP maps \cite{ouamane_sentiment_2025}, and a sentiment study finds an $R_y$ angle map beating a Pauli feature map under matched circuits \cite{ruskanda_evaluating_2025}. The attribution of any result is comparatively clean here, but at the price of a circuit that only ever sees a heavily compressed view of the input.

Hybrid neural models move the rotation downstream. The angle layer is no longer the front door but a submodule embedded in a larger classical network, so a transformer or recurrent encoder performs most of the representational work and the quantum circuit consumes a representation that has already been compressed classically. This is visible both in the front-ends, MARBERT, AraBERT, or XLM-R embeddings compressed into shallow circuits whose Pauli readouts are fused back into the classical model \cite{suleiman_hdqe_2026}, a BERT-and-GRU multimodal pipeline projected into a two-qubit circuit, a transformer generator whose tokens are injected as rotations \cite{routh_hybrid_2025}, and in the registers, which are the smallest in the angle subcorpus. The cost is the mirror image of the VQC case: because the classical encoder dominates, the independent contribution of the quantum rotation is the hardest to isolate of any family.

Recurrent models repurpose the rotation temporally rather than spatially. Here angle encoding is applied \emph{inside} a quantum recurrent cell, a QLSTM or QRNN block, so the same rotation interface is re-entered at every step of the sequence rather than once at the input \cite{stein_applying_2023,kesiku_quantum_2025,yu_application_2024,diaz-ortiz_strongly_2024}. This keeps the register deliberately narrow, a median of four qubits and never more than eight, because width is fixed by the cell while information enters through repeated application along the sequence, which brings this usage closer in spirit to data re-uploading than to a single loading step. The $R_y$-based \textsc{AngleEmbedding} is the recurring choice \cite{kesiku_quantum_2025}, and the family is the one most oriented toward token-level tasks such as part-of-speech tagging \cite{pandey_quantum_2026} alongside sequential sentiment classification.

Quantum transformers use the rotation to prepare the states that attention then compares, and this is the one place where the usual qubit discipline breaks. Because features are mapped into per-token quantum states whose overlaps play the role of attention scores, register width scales with sequence and embedding size rather than with a compressed feature count: QNet assigns one qubit per token--embedding-dimension pair and splits axes by role, with $R_x$ carrying the embedding value and $R_z$ a positional phase \cite{day_qnet_2023}, and a near-term self-attention study reaches 28 qubits, the widest register in the entire angle subcorpus \cite{he_training_2024}. Most implementations avoid this blow-up by keeping the quantum attention block small \cite{di_sipio_dawn_2022,pham_fitqt_2026}, but the family is nonetheless the only one whose structural logic pushes angle encoding beyond the near-universal $4$--$8$ qubit band.

Quantum kernels are represented by only two angle studies, and the contrast with the previous families is instructive. When angle encoding does appear in a kernel method, the rotation defines a similarity-circuit geometry rather than feeding a trainable ansatz \cite{hsu_quantum_2025,boopalamani_quantum-enhanced_2025}; but kernel methods overwhelmingly prefer purpose-built structured feature maps (Section~\ref{subsec:structured_across_constructions}), so angle encoding is a marginal, borrowed choice here rather than a native one. This is the exception that confirms the pattern: angle encoding travels everywhere, but in the one family whose encoding is itself the model, a different interface wins.

Taken together, the split shows that holding the encoding fixed does not hold the pipeline fixed. The same rotation is a complete front door in VQCs, a downstream submodule in hybrid models, a per-step cell input in recurrent models, an attention-state preparation in quantum transformers, and a borrowed similarity map in kernels, and these roles differ in exactly the dimensions that matter for evaluation: how much of the model the encoding accounts for, whether register width is set by compression or by architecture, and how cleanly the quantum contribution can be attributed. The one invariant across all five families is the heavy classical compression that precedes the circuit. This confirms the encoding--model-family correlation of Fig.~\ref{fig:encoding_model_family} from the opposite direction: even when the encoding is the same, it is the model family that determines what the encoding does, which is why angle-encoded results are not directly comparable across architectures despite sharing a label.

\subsection{Amplitude Encoding Across Model Families}
\label{subsec:amplitude_across_families}

Amplitude encoding is the natural counterpoint to angle encoding, and reading its 28 studies through the same per-family lens makes the contrast concrete. It appears in the same host families, VQCs (15), quantum kernels (5), quantum transformers (4), hybrid neural models (3), and a single recurrent model, but with an almost inverted profile, both in which families it favors and in the resources it consumes. Where angle encoding compresses text to a handful of rotation parameters over a $2$--$8$ qubit register, amplitude encoding loads a full normalized vector into the $2^{n}$ amplitudes of an $n$-qubit state, so its registers cluster near $\log_2$ of the embedding dimension rather than at a fixed small width. Table~\ref{tab:amplitude_by_family} summarizes the split, and the immediate signal is the qubit column: every family sits at a median near ten qubits, against the median of four seen for angle encoding, because the encoding is chosen precisely to avoid the aggressive projection that angle encoding requires.

\begin{table*}[t]
\caption{Amplitude encoding split by model family in the final primary corpus.}
\label{tab:amplitude_by_family}
\centering
\scriptsize
\begin{tabular}{@{}p{0.15\textwidth} p{0.06\textwidth} p{0.13\textwidth} p{0.26\textwidth} p{0.27\textwidth}@{}}
\hline
\textbf{Model family} & \textbf{Studies} & \textbf{Median qubits (range)} & \textbf{Typical classical front-end} & \textbf{Role of the amplitude state} \\
\hline
VQC & 15 & 10 (9--16) & Normalized BERT/ELECTRA/DistilBERT/SBERT embeddings & Input state prepared for a trainable ansatz \\
Quantum kernel & 5 & 10.5 (7--12) & Word/sentence embeddings, TF-IDF vectors & Overlap geometry read out by a SWAP test \\
Quantum transformer & 4 & 7.5 (4--12) & Word or contextual token embeddings & State preparation for quantum attention \\
Hybrid neural & 3 & 9 (6--12) & Multimodal or token embeddings & High-fidelity submodule after a classical encoder \\
Recurrent neural & 1 & --- & Chinese word embeddings & Per-gate state inside a decomposed QLSTM \\
\hline
\end{tabular}
\end{table*}

VQCs are again the dominant host, but the role differs from the angle case. Amplitude encoding here prepares the input state that a trainable ansatz, typically \textsc{EfficientSU2} or \textsc{RealAmplitudes}, then classifies \cite{korrapati_quantum_2025}, and the front-end is almost exclusively a dense contextual embedding normalized and padded to the nearest power of two: a 768-dimensional SBERT vector into ten qubits \cite{guarasci_applying_2025}, BERT or ELECTRA embeddings padded from 768 to 1024 amplitudes over ten qubits \cite{buonaiuto_quantum_2024}, and mBERT embeddings reused across multilingual transfer \cite{buonaiuto_multilingual_2025}. The $\log_2$ packing is stated as the point of the method rather than an implementation detail \cite{liu_hcqnsc_2026}, and where a study compares the two interfaces directly it frames amplitude encoding as the high-fidelity option and angle encoding as the cheaper fallback: one sentence-classification model reports its amplitude variant at $0.81$ accuracy against $0.71$ for a PCA-plus-angle variant of the same system \cite{guarasci_applying_2025}. This family is also markedly more recent than its angle counterpart, with every VQC amplitude study appearing from 2024 onward, which tracks the point at which transformer embeddings became the default text representation and made a dense-vector encoding worth its qubit cost.

Quantum kernels are where amplitude encoding is most at home, in contrast to the marginal, borrowed status angle encoding held there. When a normalized embedding is amplitude-encoded, the inner product of two states equals the inner product of the underlying vectors, so a SWAP test or state-overlap estimate reads out cosine-like semantic similarity directly, with no training involved \cite{cavar_old_2026,cavar_word_2025,cavar_semantic_2024}. This is why the semantic-similarity task cluster of the corpus lives almost entirely in this cell, and why kernel studies apply the same construction to fastText, GloVe, or BERT embeddings and to TF-IDF vectors feeding a quantum nearest-neighbor classifier \cite{alexander_quantum_2022,zeguendry_quantum_2024}. The fit is principled rather than incidental: the encoded state \emph{is} the geometry the kernel measures, and some studies push the density further with complex-amplitude variants that map pairs of dimensions to the real and imaginary parts of each amplitude, halving the register again \cite{cavar_old_2026}.

Quantum transformers use amplitude encoding, as they use angle encoding, to prepare the states that attention compares, but over noticeably smaller and more variable registers than in the VQC case \cite{zheng_design_2023,sharma_elevating_2025,aktar_hybrid_2026,karanath_q-optbert_2026}. Amplitude encoding is rarer here than angle encoding (four studies against nine) because preparing an arbitrary state for every token is more expensive than injecting a rotation, so the studies that do adopt it tend to encode a single pooled or sentence-level vector rather than a per-token sequence. Hybrid neural models place amplitude encoding in the same downstream-submodule role that angle encoding occupies there, but selected explicitly for fidelity: one knowledge-graph study contrasts an amplitude model chosen for high-dimensional fidelity against an angle model chosen for gate efficiency, which is the trade-off between the two families stated in a single system \cite{hendl_incorporating_2025}, while a multimodal model amplitude-encodes fused visual and textual features \cite{zheng_quantum_2025}.

The recurrent family is the sharpest point of the contrast. Amplitude encoding appears in exactly one recurrent study, a decomposed quantum LSTM over Chinese word embeddings \cite{hou_realization_2022}, against the ten recurrent studies that use angle encoding, and the asymmetry is not accidental: re-preparing an arbitrary amplitude state at every step of a sequence is prohibitively expensive, whereas re-applying a rotation inside a cell is cheap. The encoding that is native to recurrent processing is therefore the opposite of the one that is native to overlap-based comparison. Read against the angle-encoding split, amplitude encoding partitions the same families by an opposite logic: angle encoding minimizes width and pays for it in classical projection, while amplitude encoding minimizes projection and pays for it in register width and, more importantly, in state-preparation depth. That last cost is the one most often abstracted away by simulator calls, so amplitude-encoded results are the ones whose reported qubit efficiency is most sensitive to whether the price of preparing the state has actually been counted.

\subsection{Syntactic/DisCoCat Encoding Across Applications}
\label{subsec:syntactic_across_applications}

Syntactic/DisCoCat encoding cannot be split the way the previous two families were, because all 29 of its studies belong to a single model family: compositional QNLP. This is itself the first finding, and it follows from the definition given earlier, since the grammatical reduction that wires the circuit is the model, there is no separate architecture to vary. The dimension that does vary, and that plays the structural role model family played above, is the application the compositional pipeline is pointed at, because that is what determines how large and how regular the grammar-derived circuit becomes. Two properties are shared across the whole subcorpus and set it apart from angle and amplitude encoding. First, the toolchain is a near-monoculture: 22 of the 29 studies build their circuits with \textsc{lambeq}, and most of the rest use the related \textsc{DisCoPy}/\textsc{pytket} stack, in contrast to the diverse PennyLane, Qiskit, and TorchQuantum backends seen for the vector-loading families. Second, and more importantly, qubit count is not a design knob here: it is fixed by the diagram, growing with sentence length and the number of grammatical wire types rather than being chosen to fit a compression budget. Table~\ref{tab:syntactic_by_application} groups the studies by application, and the qubit column is the payoff, since it rises monotonically as the input departs from the short English sentence the formalism was built for.

\begin{table*}[t]
\caption{Syntactic/DisCoCat encoding split by application scope in the final primary corpus.}
\label{tab:syntactic_by_application}
\centering
\scriptsize
\begin{tabular}{@{}p{0.17\textwidth} p{0.06\textwidth} p{0.13\textwidth} p{0.24\textwidth} p{0.26\textwidth}@{}}
\hline
\textbf{Application scope} & \textbf{Studies} & \textbf{Median qubits (range)} & \textbf{Typical grammar/pipeline} & \textbf{Role of the grammar-derived circuit} \\
\hline
Canonical sentence classification & 12 & 4 (2--11) & \textsc{lambeq} pregroup parse to parameterized circuit & Meaning classification of grammar-labeled sentences \\
Sentiment classification & 7 & 7 (5--7) & Pregroup/\textsc{DisCoCat} with IQP or SimpleSA ansätze & Compositional sentiment over real sentences \\
Relational \& semantic tasks & 5 & 10 (single report) & Lambek/Fock, pregroup QA, \textsc{lambeq} extensions & Structure for anaphora, QA, and similarity \\
Beyond-text structured inputs & 5 & 16.5 (12--25) & Domain grammar mapped to a \textsc{DisCoCat} diagram & Grammar analog for music, SQL, and translation \\
\hline
\end{tabular}
\end{table*}

Canonical sentence classification is the heartland of the family and its smallest-circuit case, with a median of four qubits. These studies take grammar-labeled short sentences, parse them into pregroup diagrams, and classify a binary meaning, following the pipeline established by the earliest demonstrations \cite{meichanetzidis_quantum_2021,kartsaklis_lambeq_2021,lorenz_qnlp_2023}. This is also where the field's rare hardware runs are concentrated: the grammar-aware and near-term demonstrations that execute on real devices rather than in simulation sit almost entirely in this group \cite{meichanetzidis_grammar-aware_2023,lorenz_qnlp_2023}, which is possible precisely because the circuits are small. Later entries keep the same pipeline while varying peripheral choices, an optimizer comparison \cite{rajashekharaiah_sentence_2023}, an ensemble over parses \cite{bouakba_ensemble_2023}, a CCG-to-ZX route to the circuit \cite{sreedhar_towards_2025}, or natural-language-inference sentence pairs \cite{sun_quantum_2026}, but the encoding mechanism and the modest register are unchanged.

Sentiment classification applies the same pregroup machinery to sentiment-bearing sentences, and its slightly larger and tighter register, a median of seven qubits, reflects longer real-text inputs rather than curated grammar examples \cite{ganguly_quantum_2022,ruskanda_quantum_2022,peral-garcia_using_2024,peral-garcia_comparing_2024}. The most developed line here is the SimpleSA family, which fixes a concrete grammar-to-circuit recipe, an H--CNOT--$R_z$--H block with nouns left as unparameterized identity wires, and reuses it across studies \cite{ruskanda_simple_2023,ruskanda_enhancing_2024}. This group is also where the family's few departures from English appear, such as a Hindi model that parses subject--object--verb sentences into pregroup diagrams with a cup-removal step \cite{naik_extending_2026}; the fact that only two studies in the entire subcorpus leave English is a second monoculture alongside the toolchain.

Relational and semantic tasks are where compositional structure is most motivated, because the task itself is about how words relate rather than about a flat label. This group applies Lambek-calculus and Fock-space representations to pronoun and coreference resolution \cite{wazni_large_2024,wazni_towards_2023}, pregroup parsing to Chinese question answering \cite{balthasar_mansky_adapting_2023}, and \textsc{lambeq} extensions to semantic similarity and knowledge-graph relations \cite{bharti_quantum_2025}. Yet even here the compositional pipeline rarely stands alone: a similarity study, for instance, has to graft a classical embedding onto the grammar specifically to handle out-of-vocabulary words \cite{bhatuse_resolving_2025}, which shows the interface reaching the edge of what a purely grammar-derived circuit can cover.

Beyond-text structured inputs are the revealing extreme. When the compositional idea is exported to non-sentence data, the diagrams and therefore the registers become the largest in the subcorpus, a median above sixteen qubits and reaching twenty for SQL query estimation \cite{uotila_sql2circuits_2025}, thirteen for machine translation over paired source and target structures \cite{abbaszadeh_hybrid_2026}, and up to twenty-five qubits for a musical-phrase model, alongside multimodal music and video pipelines \cite{hawashin_compositional_2025,r_quantum_2023}. This is the grammar-to-circuit growth mentioned in the taxonomy, now made quantitative: the same encoding that needs two qubits for a curated sentence needs an order of magnitude more once the input is a query, a bilingual alignment, or a musical structure, because width is dictated by the diagram rather than by a compression stage that could absorb the increase.

Read together, the split shows a family whose usage varies along a single axis, distance from the short English sentence, rather than across model families. The compositional mechanism is constant and its interpretability advantage is real, but its practicality is inversely tied to that distance: circuits stay small and even reach hardware for canonical sentence classification, grow through sentiment and relational tasks, and become simulation-only and comparatively large once the input leaves ordinary sentences. Combined with the toolchain and language monocultures, this makes syntactic encoding the most internally consistent of the three established families and also the most tightly bounded, since the very grammar that gives it structure is what limits how far it can scale \cite{sreedhar_towards_2025,uotila_sql2circuits_2025}.

\subsection{Structured Feature-Map Encoding Across Constructions}
\label{subsec:structured_across_constructions}

Structured feature-map encoding is the smallest of the four families addressed here, with nine studies, and it repeats the concentration seen for syntactic encoding: six of the nine are quantum-kernel methods, with only single instances in VQC, hybrid neural, and quantum-transformer models. Because model family again barely varies, the dimension that plays its role is the feature-map construction itself, which is the natural axis for this family in any case: unlike angle and amplitude encoding, which load feature \emph{values}, a structured map builds feature \emph{interactions} into the circuit through phase rotations and entangling gates, so the choice of map is a modeling assumption about which interactions matter. Two properties frame the split. Like angle encoding, the family depends on aggressive classical compression before the circuit, since the high-dimensional text representation is reduced by PCA or a projection layer first; unlike angle encoding, register width then tracks the number of retained features at roughly one qubit each, which is why the reported registers are comparatively wide, with a median of $10.5$ qubits and a range of $5$ to $20$. Table~\ref{tab:structured_by_construction} groups the studies by the map they use.

\begin{table*}[t]
\caption{Structured feature-map encoding split by map construction in the final primary corpus.}
\label{tab:structured_by_construction}
\centering
\scriptsize
\begin{tabular}{@{}p{0.17\textwidth} p{0.06\textwidth} p{0.13\textwidth} p{0.24\textwidth} p{0.26\textwidth}@{}}
\hline
\textbf{Map construction} & \textbf{Studies} & \textbf{Median qubits (range)} & \textbf{Interactions encoded} & \textbf{Host families and role} \\
\hline
ZZ / IQP maps & 4 & 10 (5--20) & Pairwise, second-order feature products via entangling phase gates & Mostly kernel geometry; one hybrid encoding layer \\
General Pauli maps & 3 & 11 (6--14) & Higher-order and non-$Z$ Pauli-string interactions & Kernel and VQC; gate combination chosen by search \\
Unnamed structured maps & 2 & not reported & Unspecified feature interactions & QSVM branch; transformer feature stage \\
\hline
\end{tabular}
\end{table*}

ZZ and IQP maps are the default construction, accounting for four of the nine studies. These are second-order maps that entangle pairs of features through data-dependent phase gates, and they are most at home in quantum kernels, where the encoded state defines the similarity geometry that the kernel measures: a ZZFeatureMap QSVM over a twenty-qubit register \cite{widdows_near-term_2024}, an IQP feature map applied to reduced AraBERT embeddings \cite{djemmal_arabert-qc_2025}, and IQP/ZZ maps feeding projected quantum kernels \cite{keith-norambuena_quantum_2026}. The same construction also appears once outside a kernel, as the primary quantum encoding layer of a hybrid model whose BERT/T5 embeddings are projected into a five-qubit ZZ map \cite{baik_neural_2025}, which illustrates the family's borderline status: even its standard map is used both as a kernel definition and as a swappable layer inside a trainable network.

General Pauli maps, in three studies, are the more expressive and more experimental construction. These use higher-order or non-$Z$ Pauli strings and treat the exact gate content as a design choice to be tuned: a question-answering VQC selects a specific $X$/$Y$/$ZZ$ combination after a small search \cite{katyayan_supervised_2023}, and a quantum-enhanced SVM adopts a Pauli-$Y$/$YYY$ feature map with a rotation factor over fourteen qubits produced by PCA \cite{ruskanda_quantum-enhanced_2023}, alongside an earlier Pauli-map sentiment kernel \cite{mohammad_masum_hybrid_2023}. This is exactly where the family's central risk becomes visible: because the map fixes which feature interactions are represented, performance depends on whether that hand-chosen structure happens to match the linguistic task, and the map is a strong assumption rather than a learned one. The remaining two studies report only an unnamed ``structured feature map'', one as a QSVM branch compared against a VQC \cite{sunkavalli_experimentation_2024} and one as a feature stage inside an aviation-text transformer \cite{jagatheesan_quantum_2026}, a reporting gap that makes them hard to place against the named constructions.

Taken together, structured feature maps are the most-benchmarked yet least-dominant of the established families. Because the map is usually presented as one candidate encoding among several, the studies that use it tend to compare it directly with the vector-loading families, and those comparisons do not favor it: the same research group behind the Pauli-$Y$ kernel above later reports, in the controlled study cited in the angle-encoding analysis, that a simple $R_y$ angle map outperforms a Pauli feature map under matched circuits \cite{ruskanda_evaluating_2025}. The concentration in kernels is therefore the coherent core of the family, the one setting where the map is the model rather than a replaceable layer, while the scattered VQC, hybrid, and transformer instances are experiments in borrowing a designed geometry into a trainable pipeline. This leaves the family with a narrow but principled niche, testing whether an engineered similarity space helps for text, subject to a double dependency the other families do not share: its usefulness rests on both the feature-map design and the compression stage that precedes it.

\subsection{Emerging and Specialized Interfaces}
\label{subsec:emerging_specialized_interfaces}

The residual group of eight studies is best read not as one family but as three distinct mechanisms, and unlike the four established families it cannot be split by model family for a different reason: it is already almost mono-cultural in architecture. Seven of the eight studies are variational or recurrent circuits, with a single hybrid-neural case and no kernel, compositional, or transformer instances at all. What varies, and what organizes this subsection, is the specialized mechanism itself, because each one is best understood as a targeted patch for a specific limitation of the mainstream encodings analyzed above rather than as a general-purpose alternative to them. Table~\ref{tab:specialized_by_mechanism} groups the studies by mechanism and names the trade each one makes.

\begin{table*}[t]
\caption{Specialized encoding interfaces split by mechanism in the final primary corpus.}
\label{tab:specialized_by_mechanism}
\centering
\scriptsize
\begin{tabular}{@{}p{0.16\textwidth} p{0.06\textwidth} p{0.12\textwidth} p{0.28\textwidth} p{0.24\textwidth}@{}}
\hline
\textbf{Mechanism} & \textbf{Studies} & \textbf{Median qubits (range)} & \textbf{Core mechanism} & \textbf{Host families and trade} \\
\hline
Data re-uploading & 4 & 7 (6--18) & Same features re-injected across circuit depth, alternating with trainable gates & VQC, recurrent; trades qubit width for depth \\
Hybrid embedding & 3 & 6 (4--8) & A classical encoder emits a compact latent vector consumed by a quantum layer & VQC, hybrid; trades quantum role for classical preprocessing \\
Hamiltonian/observable & 1 & 9 & Data mapped to Pauli-string Hamiltonians or observables, with no state preparation & VQC; trades state encoding for observable encoding \\
\hline
\end{tabular}
\end{table*}

Data re-uploading is the response to the width problem that constrains angle and amplitude encoding. Instead of loading features once and then applying only trainable gates, it re-injects the same feature stream across the depth of the circuit, alternating data-dependent rotations with trainable transformations, so a small register can carry more information than a single loading step would allow. The four studies make this trade explicitly: a bricklayer data-reuploading classifier over six qubits \cite{hamilton_characterizing_2023}, a lexical-feature circuit that trades qubit width for repeated injection \cite{silver_lexiql_2024}, a recurrent variant that re-uploads code-token embeddings across recurrent quantum layers for vulnerability detection \cite{song_recurrent_2024}, and a sentence-transformer pipeline that spreads a large embedding across depth over eighteen qubits \cite{kim_quantum_2025}. The register range, from six to eighteen qubits, reflects exactly this: width is no longer minimized to a fixed small band as in angle encoding, but neither does it grow to $\log_2$ scale as in amplitude encoding, because the burden has been shifted into circuit depth instead.

Hybrid embedding is the response to the high-dimensionality problem, and it is a system-level interface rather than a gate-level rule. A classical encoder, an autoencoder or a multimodal network, first compresses the text, and sometimes other modalities, into a compact latent vector, which a small quantum layer then consumes. This is why these are the smallest registers in the specialized group, a median of six qubits: an eight-dimensional latent from a text-and-image emotion model mapped to eight qubits \cite{castrignano_toward_2025}, a minimal intermediate representation of code-token embeddings over four qubits \cite{akter_automated_2023}, and a multimodal misinformation pipeline that enriches GloVe and visual features through a quantum fusion layer \cite{bikku_quantum-enhanced_2025}. The recurring theme, and the recurring weakness, is that all three are multimodal or code-analysis systems in which the classical encoder does most of the representational work, so the independent contribution of the small quantum tail is the hardest of any interface to isolate, which is the same attribution problem seen for hybrid-neural angle encoding, now taken to its limit.

Hamiltonian and observable encoding is the most conceptually distinct of the three, and it is the response to the state-preparation problem that undermines amplitude encoding's qubit savings. A single study maps text-derived data to Pauli-string Hamiltonians used as observables rather than preparing a quantum state at all, which sidesteps the state-preparation circuit entirely while keeping a logarithmic qubit count, demonstrated over nine qubits for sentiment classification \cite{tiblias_efficient_2026}. This changes the object that carries the linguistic information, from a prepared state to a measured observable, and so directly targets the cost that Section~\ref{subsec:amplitude_across_families} identified as amplitude encoding's hidden liability. With only one study, however, it is an idea rather than an established interface.

Taken together, the specialized interfaces confirm the reading of the established families from the outside in: each is defined by the mainstream limitation it tries to escape, data re-uploading trading width for depth, hybrid embedding trading quantum responsibility for classical preprocessing, and Hamiltonian encoding trading state preparation for observables. Their shared confinement to variational and recurrent circuits shows that this experimentation is happening entirely within the trainable-circuit world, not in the kernel or compositional families whose encodings are structurally fixed. The evidence, however, is thinner and less comparable than for any established family, because these studies differ not only in encoding but in where the quantum component sits and how much work precedes it, so the natural questions, whether depth is worth reduced width, whether the quantum layer adds value beyond the classical encoder, and whether observable encoding beats state preparation, remain open on the current corpus.

\subsection{NLP Tasks and Encoding Choice}
\label{subsec:tasks_encodings}

\Figure[t!](topskip=0pt, botskip=0pt, midskip=0pt)[width=\textwidth]{figures/encoding_by_task_family.pdf}{Encoding family by NLP task family in the final primary corpus. Studies evaluating tasks from several families are counted once per family.\label{fig:encoding_task}}

\Figure[t!](topskip=0pt, botskip=0pt, midskip=0pt)[width=\textwidth]{figures/task_family_by_year.pdf}{NLP task families addressed per publication year. Each bar segment counts studies evaluating at least one task from the family.\label{fig:task_year}}

To characterize what the encodings are actually used \emph{for}, the 31 raw task labels extracted from the corpus have been grouped into seven task families: sentiment classification, other classification (topic, spam, misinformation, acceptability, and similar detection tasks), similarity and disambiguation, sequence labeling, question answering and retrieval, generation and language modeling, and a residual group. In this scheme, a study that evaluates tasks from several families has been counted once in each of them.

The resulting distribution is heavily skewed toward classification. In fact, ninety of the 125 studies (72\%) evaluate at least one classification task, and 80 of them (64\%) evaluate nothing else: sentiment classification appears in 47 studies and other classification variants in 53. By comparison, the task families that define much of contemporary NLP remain marginal: similarity and disambiguation appear in 13 studies, generation and language modeling in 10, question answering and retrieval in 6, and sequence labeling in 6, while machine translation, dialogue, and code generation are each represented by a single study. QNLP, as currently practiced, is therefore largely a study of small-scale text classification.

Fig.~\ref{fig:encoding_task} crosses encoding family with task family, and it shows that the encodings are not distributed uniformly over this task space. Angle encoding, despite being the broadest family, is nearly absent from similarity work, where it appears in a single study; amplitude encoding, by contrast, carries that task family almost alone (6 of 13 studies), which is consistent with its natural fit to overlap- and inner-product-based comparison \cite{cavar_old_2026,cavar_word_2025}. Syntactic/DisCoCat encoding, in turn, is not confined to toy sentence classification: it also reaches disambiguation and similarity (5 studies) and question answering (2), where grammatical structure is plausibly informative. Sequence labeling, on the other hand, is almost exclusively angle-encoded (5 of 6 studies), which reflects the token-by-token injection pattern of quantum recurrent and transformer-style models. In other words, encoding choice partially tracks task structure, but the dominant driver remains the model family in which the encoding is embedded.

Fig.~\ref{fig:task_year}, finally, shows the temporal development of this task space. Output grows steeply, from 2 studies in 2021 to 41 in 2025, and yet the task mix barely changes: from 2022 onward, classification accounts for between 63\% and 81\% of each year's studies, with no visible diversification trend. The absolute number of generation and language-modeling studies has risen slightly since 2024, but their share has remained flat. In short, the field is growing without broadening.

\subsection{Datasets and Cross-Study Comparability}
\label{subsec:datasets_comparability}

A second lens on the corpus is provided by the datasets on which the encodings are evaluated. After naming variants have been normalized (e.g., IMDB/IMDb, or the several designations of the \textsc{lambeq} example sets), the 118 studies that name at least one dataset are found to use 141 distinct datasets, whereas 7 studies do not identify their data at all. Of these 141 datasets, 122 are used by exactly one study, only 19 are shared by two or more studies, and only 7 reach four or more: the \textsc{lambeq} MC (12 studies) and RP (9) example sets \cite{meichanetzidis_quantum_2021,lorenz_qnlp_2023}, IMDb (12), Amazon reviews (8), Yelp (7), SST-2 (5), and the restaurant-domain sentence set introduced with SimpleSA \cite{ruskanda_simple_2023} (5). In total, only 50 of the 125 studies share at least one dataset with any other study in the corpus.

Two main consequences follow from this fragmentation. First, the most widely reused ``benchmarks'' in QNLP are the MC and RP sets, that is, artificial corpora of roughly 130 short sentences over vocabularies of a few dozen words, two to three orders of magnitude smaller than standard NLP benchmarks; and where established sentiment benchmarks such as IMDb or SST-2 are used, they are typically subsampled to a few hundred or a few thousand examples in order to fit simulation budgets. Second, cross-study performance comparison is possible only inside the handful of shared-dataset cells. In fact, the MC cell is the most instructive of these: it spans syntactic, angle, amplitude, and data re-uploading encodings between 2021 and 2026, with qubit budgets ranging from 2 to 28, and with reported accuracies approaching or reaching 100\%. Such near-ceiling results on a 130-sentence corpus, however, say little about relative encoding quality, and outside these cells the studies differ simultaneously in dataset, subsampling, metric, and baseline, so that no meta-analytic ranking of encoding families is currently supportable. For this reason, the per-study extraction table released with this survey is intended to make such comparisons possible as shared evaluation practice matures.

\subsection{Dataset Scale and Qubit Budgets}
\label{subsec:scale_qubits}

\Figure[t!](topskip=0pt, botskip=0pt, midskip=0pt)[width=\textwidth]{figures/dataset_size_vs_qubits.pdf}{Dataset size versus qubit count for the 37 studies reporting both, colored by encoding family. The two quantities are uncorrelated ($r \approx -0.07$ between $\log_{10}$ dataset size and qubit count).\label{fig:size_qubits}}

If the encoding step mediated between data scale and quantum resources, then larger datasets should tend to engage larger qubit registers. Fig.~\ref{fig:size_qubits}, however, shows that they do not. Across the 37 studies that report both quantities, dataset sizes are found to span four orders of magnitude (from 30 to 500{,}000 examples), whereas qubit counts remain confined to a narrow band (from 2 to 28, with a median of 8), and the correlation between the logarithm of dataset size and qubit count is effectively zero ($r \approx -0.07$). This decoupling is visible even at the extremes: the largest dataset in the corpus, a 500{,}000-example sentiment collection, is processed with 5 qubits \cite{rawat_mental_2026}, whereas the largest register, of 28 qubits, is applied to the 130-sentence MC set \cite{he_training_2024}.

Taken together, this decoupling can be read as the clearest corpus-level signature of the encoding bottleneck. Qubit width has in fact been determined by the classical compression stage, i.e., by how many features survive projection, rather than by the linguistic problem itself, so that growth in data scale is absorbed entirely before the quantum circuit is reached. The same analysis also exposes a reporting gap that limits any resource analysis: only 80 of the 125 studies report a usable qubit count, 52 report a dataset size, and just 37 report both. For this reason, reporting both quantities, together with the dimensionality of the classical representation at the encoding interface, should become standard practice in future QNLP work.

\bibliography{references}
\bibliographystyle{IEEEtran}

\EOD

\end{document}
